% Embedded Systems lab report -- compile with pdflatex (twice, for the ToC)
\documentclass[11pt,a4paper]{article}

\usepackage[margin=2.5cm]{geometry}
\usepackage{lmodern}      % scalable Type1 fonts for T1; without this,
\usepackage[T1]{fontenc}  % MiKTeX builds slow, fuzzy METAFONT bitmaps
\usepackage[utf8]{inputenc}
\usepackage{graphicx}
\usepackage{booktabs}
\usepackage{xcolor}
\usepackage{listings}
\usepackage[hidelinks]{hyperref}

% --- Listing styles --------------------------------------------------------
\definecolor{codegray}{gray}{0.96}
\definecolor{codegreen}{rgb}{0.0,0.5,0.0}
\definecolor{codeblue}{rgb}{0.0,0.0,0.6}

\lstdefinestyle{pystyle}{
  language=Python,
  backgroundcolor=\color{codegray},
  basicstyle=\ttfamily\footnotesize,
  keywordstyle=\color{codeblue}\bfseries,
  commentstyle=\color{codegreen}\itshape,
  stringstyle=\color{red!70!black},
  numbers=left,
  numberstyle=\tiny\color{gray},
  numbersep=8pt,
  breaklines=true,
  frame=single,
  framerule=0pt,
  showstringspaces=false,
  tabsize=4,
  captionpos=b,
}
\lstdefinestyle{shellstyle}{
  language=bash,
  backgroundcolor=\color{codegray},
  basicstyle=\ttfamily\footnotesize,
  breaklines=true,
  frame=single,
  framerule=0pt,
  showstringspaces=false,
  captionpos=b,
}
\lstset{style=pystyle}

% --- Title ----------------------------------------------------------------
\newcommand{\studentid}{R12942157}

\title{Embedded Systems Laboratory Report\\
       \large BLE Central Programming with a Raspberry Pi\\
       \normalsize Basic Problem and Option Problem 1}
\author{YU-CHEN, CHU\\
        \normalsize Student ID: \studentid}
\date{\today\\[2ex]
      \normalsize Source code:
      \url{https://github.com/oldpig123/embedded_system/tree/master/lab3}}

\begin{document}
\maketitle
\tableofcontents
\newpage

% =========================================================================
\section{Introduction}

This exercise is about the client side of Bluetooth Low Energy (BLE). A
Raspberry Pi acts as a BLE \emph{central}, which in GATT terms is the
\emph{client}, and connects to a \emph{peripheral}, the GATT \emph{server}.
The task is to change the value of a Client Characteristic Configuration
Descriptor (CCCD) on the server from the Python program running on the Pi.

The CCCD (UUID \texttt{0x2902}) is the mechanism by which a client controls
which server-initiated packets it receives for a given characteristic. A
server may send a notification only if the client has written \texttt{0x0001}
to the CCCD, and an indication only if it has written \texttt{0x0002}. The
value is a 16-bit little-endian integer, so the byte sequence actually sent
over the air for \texttt{0x0002} is \texttt{02 00}. This byte order is an easy
mistake to make and is called out again in the implementation section.

The report covers the basic problem (a Python central that writes
\texttt{0x0002} to a CCCD) and the optional problem (building the C library
GATTLIB on the Pi and trying its examples).

% =========================================================================
\section{Environment}

Two machines are involved. The Raspberry Pi is the device under test and runs
the central. A Windows laptop stands in for the phone and runs a small GATT
server (see Section~\ref{sec:deviation} for why).

\begin{table}[h]
\centering
\caption{Software and hardware used.}
\begin{tabular}{@{}lll@{}}
\toprule
 & Raspberry Pi (central) & Windows laptop (peripheral) \\
\midrule
OS            & Debian 13 (``trixie'')           & Windows 11 \\
Bluetooth     & Pi onboard adapter (\texttt{hci0}) & Qualcomm FastConnect 7800 \\
BLE stack     & BlueZ 5.82                       & Windows WinRT GATT API \\
Language      & Python 3.13, C (gcc 14.2)        & Python 3.14 (managed by \texttt{uv}) \\
Libraries     & \texttt{bluepy} 1.3.0, GATTLIB    & \texttt{winrt-*} 3.2.1 \\
\bottomrule
\end{tabular}
\end{table}

\subsection{Raspberry Pi setup}

The lecture notes install \texttt{bluepy} with \texttt{sudo pip install}. On
Debian 13 this is refused with \texttt{error: externally-managed-environment}
(PEP~668): the distribution prevents \texttt{pip} from writing into the system
Python so that it cannot break OS packages. The fix is a virtual environment:

\begin{lstlisting}[style=shellstyle,caption={Installing bluepy on Debian 13.}]
sudo apt install python3-pip libglib2.0-dev
python3 -m venv ~/blenv
~/blenv/bin/pip install bluepy
sudo ~/blenv/bin/python central.py      # root is needed for BLE scanning
\end{lstlisting}

Because the program must run as root, \texttt{sudo} has to be given the
virtual environment's interpreter explicitly. A plain \texttt{sudo python}
would use the system Python, in which \texttt{bluepy} is not installed.

% =========================================================================
\section{Basic Problem: CCCD Write from a Python Central}
\label{sec:basic}

\subsection{Design}

The program follows the scan--connect--write sequence of the lecture notes:

\begin{enumerate}
  \item Scan for 8 seconds and select the peripheral that advertises service
        \texttt{0xFFF0}.
  \item Connect to it.
  \item Find characteristic \texttt{0xFFF1} and, from it, its CCCD
        (\texttt{0x2902}).
  \item Write the two bytes \texttt{02 00} (the value \texttt{0x0002}) with a
        write \emph{request}, so that the server's response, or its error, is
        observed.
  \item Read the descriptor back and print it, then disconnect in a
        \texttt{finally} block so that the link is released even if a step
        fails.
\end{enumerate}

\subsection{Selecting the peripheral}

An unfiltered scan in a normal room returned about 70 devices, so the target
must be identified by content rather than by position in the list. The
peripheral's address is not usable for this: it is a \emph{random} address and
\emph{changes} from one advertising session to the next (the same laptop
appeared under several different addresses across runs). The program therefore
matches on the advertised service UUID and then uses whatever address that
entry carries, passing the address type \texttt{'random'} to
\texttt{Peripheral}. Without the address type, \texttt{bluepy} fails with
``connection refused (111)'', as noted in the lecture slides.

Two details of the advertisement matter. First, \texttt{getValueText()}
returns \texttt{None} for entries that lack the requested field, which is true
of most devices in a busy room, so the \texttt{None} case must be skipped
before the substring test. Second, the target UUID was found under AD type~3
(complete list of 16-bit service UUIDs), not AD type~7 (128-bit), because the
\texttt{0xFFF0} UUID is a 16-bit short form of the Bluetooth base UUID. A first
attempt that looked only at type~7 found nothing even though the device was
being received.

\subsection{Implementation}

\lstinputlisting[caption={\texttt{central.py}: BLE central running on the Raspberry Pi.}]{../central.py}

The CCCD is located through the \emph{characteristic} object
(\texttt{char.getDescriptors(forUUID=0x2902)}), not through the peripheral:
\texttt{Peripheral.getDescriptors} takes a handle range and has no
\texttt{forUUID} argument. The value is written as \verb|b"\x02\x00"|; the bytes \texttt{00 02} would instead set the value
to \texttt{0x0200}.

\subsection{The peripheral used for testing}
\label{sec:deviation}

The assignment suggests an Android phone running BLE Tool or a BLE scanner
app. Installing an app on the phone was not desired, so the peripheral role was
instead played by the Windows laptop. Its Bluetooth adapter was confirmed to
support the BLE peripheral role (\texttt{IsPeripheralRoleSupported = True}
through the WinRT \texttt{BluetoothAdapter} API). A short Python program,
\texttt{server.py}, advertises service \texttt{0xFFF0} containing one
\emph{indicate-only} characteristic \texttt{0xFFF1}, and prints a timestamped
line whenever a client's subscription changes.

\lstinputlisting[caption={\texttt{server.py}: GATT server on the laptop, standing in for the phone.}]{../server.py}

The characteristic is deliberately indicate-only, so that the only valid CCCD
value is \texttt{0x0002}.

The choice of library was not the first one. The popular \texttt{bless}
package was tried first, but version 0.2.6 does not work on current Windows: it
imports the obsolete \texttt{bleak\_winrt} module, requires a
\texttt{pysetupdi} package that is not on PyPI, and, to set the device name, edits the adapter's
registry entry. The server was therefore written directly
against the WinRT \texttt{GattServiceProvider} classes, which needs only a few
calls.

\paragraph{A limitation of this peripheral.} On Windows the operating system
owns the CCCD. Unlike the Android BLE Tool, whose server log shows the raw
descriptor write, a Windows server program can observe only the
\texttt{SubscribedClientsChanged} event, i.e.\ how many clients are
subscribed, not the \texttt{0x0002} value itself. The evidence for the write
is therefore the combination of two independent observations: the Pi reads the
descriptor back as \texttt{0x0002}, and the laptop's server log records the
subscription. The indicate-only property is what ties a subscription to the
value \texttt{0x0002}.

\subsection{Results}

With \texttt{server.py} running on the laptop, \texttt{central.py} found the
device, wrote the descriptor, and read back the value \texttt{0x0002}. The
laptop's log showed a subscription at the moment of the write:

\begin{lstlisting}[style=shellstyle,caption={Laptop log during one run of the central.}]
Advertising service 0000fff0-0000-1000-8000-00805f9b34fb. Ctrl+C to stop.
[15:48:11] CCCD changed on 0000fff1-...: 1 subscribed client(s)
[15:48:13] CCCD changed on 0000fff1-...: 0 subscribed client(s)
\end{lstlisting}

The second line is \emph{not} a second write. It is the cleanup that occurs
when the Pi disconnects in the \texttt{finally} block: a CCCD subscription does
not survive the connection, so the server sees the subscriber count return to
zero. The two timestamps therefore show subscribe at the write and unsubscribe
at the disconnect.

% =========================================================================
\section{Problems Encountered}

Most of the time in this lab was spent not on the CCCD logic itself but on
getting a working BLE environment. The causes are listed because each one
produces a misleading symptom.

\begin{table}[h]
\centering
\caption{Problems, their symptoms and causes.}
\begin{tabular}{@{}p{3.2cm}p{4.8cm}p{5.8cm}@{}}
\toprule
Problem & Symptom & Cause and fix \\
\midrule
\texttt{pip install} refused &
  \texttt{externally-managed-environment} &
  Debian 13 enforces PEP~668. Use a venv and run it with
  \texttt{sudo} by absolute path. \\
Scan finds nothing, \texttt{scanend ... Rejected} &
  \texttt{bluepy} raised \texttt{BTLEManagementError} after the scan; empty
  device list; \texttt{bluetoothctl} said \texttt{NotReady} &
  Bluetooth was \emph{soft-blocked} (\texttt{rfkill list}). Fixed with
  \texttt{rfkill unblock bluetooth} and \texttt{hciconfig hci0 up}. The error
  message pointed at \texttt{bluepy}, but the library was not at fault. \\
\texttt{bless} fails on Windows &
  \texttt{No module named bleak\_winrt}, then \texttt{pysetupdi} &
  Unmaintained dependencies. Replaced by direct WinRT calls. \\
Laptop invisible to the Pi &
  Filter on AD type~7 found nothing &
  The UUID is in AD type~3. A raw dump of all scan fields exposed it. \\
Address kept changing &
  A stored MAC no longer matched &
  Random advertising addresses rotate; select by service UUID instead. \\
\bottomrule
\end{tabular}
\end{table}

The soft-blocked adapter was the most instructive case. Two unrelated tools,
\texttt{bluepy} and \texttt{bluetoothctl}, failed with different messages, and
only the lower-level \texttt{rfkill} state explained both. When a stack of
libraries reports an error, the fault should be looked for at the lowest layer
first.

% =========================================================================
\section{Option Problem 1: BLE in C with GATTLIB}

\subsection{Building GATTLIB}

GATTLIB was cloned from \url{https://github.com/labapart/gattlib} and built
with CMake on the Pi. CMake selected the D-Bus backend for BlueZ~5.82. The
build initially failed at configuration time, in the examples only, with
\texttt{None of the required 'libpcre' found}.

The cause is that the examples' \texttt{CMakeLists.txt} files request the
PCRE1 library (\texttt{libpcre}), while Debian 13 no longer ships
\texttt{libpcre3-dev} (\texttt{apt-cache policy} reported no candidate).
PCRE2 (\texttt{libpcre2-dev}) was already installed as a dependency of GLib.
Before changing anything, a search of the C sources showed that none of the
GATTLIB code actually calls the PCRE API; only the CMake files link it. The
one-word change in eight files, from \texttt{libpcre} to \texttt{libpcre2-8},
was therefore enough:

\begin{lstlisting}[style=shellstyle,caption={Build steps, including the PCRE workaround.}]
sudo apt install git cmake build-essential libbluetooth-dev libreadline-dev
git clone https://github.com/labapart/gattlib.git && cd gattlib
grep -rl "REQUIRED libpcre)" --include=CMakeLists.txt . \
  | xargs sed -i 's/REQUIRED libpcre)/REQUIRED libpcre2-8)/'
mkdir build && cd build && cmake .. && make
sudo make install
\end{lstlisting}

The build and installation completed. The examples were run from the
\texttt{build} directory. The compiler emitted warnings (transposed
\texttt{calloc} arguments, an unused variable) from the library's own code
under gcc~14, which do not affect the result.

\subsection{Trying the examples}

\paragraph{\texttt{ble\_scan}.} This example scans and then attempts to
connect to \emph{every} device it found, listing each one's GATT services. In
a room with several dozen random-address devices most connections fail, for
example with \texttt{Timeout was reached} or
\texttt{le-connection-abort-by-local}. These failures are expected rather than
a fault in the setup: the devices are other people's phones, trackers and
peripherals that do not accept a connection from a stranger. One device did
accept and its services (\texttt{0x1800}, \texttt{0x1801}, \texttt{0x180A},
\texttt{0x180F} and several vendor-specific UUIDs) were listed, which
demonstrates that scanning, connecting and service discovery all work. The
program also terminated with a segmentation fault when it attempted to connect
to a network printer. This crash was observed but not investigated; it occurs
in the example's failed-connection path, not in the code written for this lab.

When the output was piped to \texttt{tee} for logging, the log initially came
out empty. This is standard C buffering: when standard output is a pipe it is
block-buffered, and the process was interrupted with Ctrl+C before the buffer
was flushed. Running it through \texttt{stdbuf -oL} (line-buffered) and
redirecting standard error into the same pipe fixed it.

\paragraph{\texttt{discover}.} Aimed at the laptop running
\texttt{server.py}, \texttt{discover} listed the laptop's complete GATT
table, which includes many built-in Windows services, and ended with the
characteristic written by \texttt{server.py}:

\begin{lstlisting}[style=shellstyle,caption={The last line of \texttt{discover}'s output.}]
characteristic[36] properties:20 value_handle:0068 uuid:0xfff1
\end{lstlisting}

The property byte \texttt{0x20} is the GATT \emph{Indicate} bit, which
matches how \texttt{server.py} defined the characteristic, so GATTLIB read the
true table from the laptop. The \texttt{read\_write} and \texttt{notification}
examples need a readable characteristic and a Battery service respectively, and
neither exists on the test peripheral.

% =========================================================================
\section{Conclusion}

A Python BLE central was written that scans, selects its peripheral by
advertised service, connects, writes \texttt{0x0002} to a CCCD with a write
request, reads the value back and disconnects cleanly. It was tested against a
purpose-written GATT server on a Windows laptop instead of an Android phone.
The principal limitation of that choice is that Windows hides the raw CCCD
value from the server application, which is why the evidence combines the
Pi's read-back with the server's subscription event.

In the optional part, GATTLIB was built on a Debian~13 system for which its
documented dependencies are partly out of date, and two of its examples were
used to scan and to enumerate the services of a device under test.

Beyond the CCCD itself, the main lesson was how much of BLE programming is
about the environment: soft-blocked radios, address types and rotating
addresses, byte order, and libraries that have not kept pace with the
operating system. Each of those produced a symptom far from its cause, and each
was solved by testing the layer below the one that reported the error.

\end{document}
